\SourceLeaf{024}
$Z$ rationaliter determinabilem fore asserit clar. de F., et proin realem. Hinc sequitur, radices aequationis $u'u'+u''u'+M''=0$ sub forma $p+q\BareRoot{-1}$ contentas fore; eaedem vero manifesto aequationi $U'=0$ satisfacient: quare dabitur valor aliquis ipsius $u'$ sub forma $p+q\BareRoot{-1}$ contentus. Iam coëfficiens $M'$ (eodem modo ut ante) rationaliter per $u'$ et coëfficientes in $Z$ determinari potest, adeoque etiam sub forma $p+q\BareRoot{-1}$ contentus erit; quare aequationis $uu+u'u+M'$ radices sub eadem forma contentae erunt, simul vero aequationi $U=0$ satisfacient, i. e. aequatio haec habebit radicem sub forma $p+q\BareRoot{-1}$ contentam. Denique hinc simili ratione sequitur, etiam $M$ sub eadem forma contineri, nec non radicem aequationis $zz+uz+M=0$, quae manifesto etiam aequationi propositae $Z=0$ satisfaciet. Quamobrem quaevis aequatio ad minimum unam radicem formae $p+q\BareRoot{-1}$ habebit.

\Article{11}
Obiectiones 1, 2, 3, quas contra \textsc{Euleri} demonstrationem primam feci (art. 8), eandem vim contra hanc methodum habent, ea tamen differentia, ut obiectio secunda, cui \textsc{Euleri} demonstratio tantummodo in quibusdam casibus specialibus obnoxia erat, praesentem in omnibus casibus attingere debeat. Scilicet a priori demonstrari potest, etiamsi formula detur, quae coëfficientem $M'$ rationaliter per $u'$ et coëfficientes in $Z$ exprimat, hanc pro pluribus valoribus ipsius $u'$ necessario indeterminatam fieri debere; similiterque formulam, quae coëfficientem $M''$ per $u''$ exhibeat, indeterminatam fieri pro quibusdam valoribus ipsius $u''$ etc. Hoc luculentissime perspicitur, si aequationem quarti gradus pro exemplo assumimus. Ponamus itaque $m=4$, sintque radices aequationis $Z=0$, hae $\alpha$, $\beta$, $\gamma$, $\delta$. Tum patet, aequationem $U=0$ fore sexti gradus ipsiusque radices $-(\alpha+\beta)$, $-(\alpha+\gamma)$, $-(\alpha+\delta)$, $-(\beta+\gamma)$, $-(\beta+\delta)$, $-(\gamma+\delta)$. Aequatio $U'=0$ autem erit decimi quinti gradus, et valores ipsius $u'$ hi
\[
\begin{gathered}
2\alpha+\beta+\gamma,\quad 2\alpha+\beta+\delta,\quad 2\alpha+\gamma+\delta,\\
2\beta+\alpha+\gamma,\quad 2\beta+\alpha+\delta,\quad 2\beta+\gamma+\delta,\\
2\gamma+\alpha+\beta,\quad 2\gamma+\alpha+\delta,\quad 2\gamma+\beta+\delta,\\
2\delta+\alpha+\beta,\quad 2\delta+\alpha+\gamma,\quad 2\delta+\beta+\gamma,\\
\alpha+\beta+\gamma+\delta,\quad\alpha+\beta+\gamma+\delta,\quad\alpha+\beta+\gamma+\delta
\end{gathered}
\]
Iam in hac aequatione, quippe cuius gradus est impar, subsistendum erit, habebitque ea revera radicem realem $\alpha+\beta+\gamma+\delta$ (quae primo coëfficienti in $Z$ mutato signo aequalis adeoque non modo realis sed etiam rationalis erit. si coëfficien\SourceLeaf{025}tes in $Z$ sunt rationales). Sed nullo negotio perspici potest, si formula detur, quae valorem ipsius $M'$ per valorem respondentem ipsius $u'$ rationaliter exhibeat, hanc necessario pro $u'=\alpha+\beta+\gamma+\delta$ indeterminatam fieri. Hic enim valor \textit{ter} erit radix aequationis $U'=0$, respondebuntque ipsi tres valores ipsius $M'$, puta $(\alpha+\beta)(\gamma+\delta)$, $(\alpha+\gamma)(\beta+\delta)$ et $(\alpha+\delta)(\beta+\gamma)$, qui omnes irrationales esse possunt. Manifesto autem formula rationalis neque valorem irrationalem ipsius $M'$ in hoc casu producere posset, neque tres valores diversos. Ex hoc specimine satis colligi potest, methodum clar. \textsc{de Foncenexii} neutiquam esse satisfacientem, sed si ab omni parte completa reddi debeat, multo profundius in theoriam eliminationis inquiri oportere.

\Article{12}
Denique ill. \textsc{La Grange} de theoremate nostro egit in comm. \textit{Sur la forme des racines imaginaires des équations, Nouv. Mém. de l'Acad. de Berlin} 1772, p. 222 sqq. Magnus hic geometra imprimis operam dedit, defectus in \textsc{Euleri} demonstratione prima supplere et revera praesertim ea, quae supra (art. 8) obiectionem secundam et quartam constituunt, tam profunde perscrutatus est. ut nihil amplius desiderandum restet, nisi forsan in disquisitione anteriori super theoria eliminationis (cui investigatio haec tota innititur) quaedam dubia superesse videantur. --- Attamen obiectionem tertiam omnino non attigit, quin etiam tota disquisitio superstructa est suppositioni, quamvis aequationem $m^{\mathrm{ti}}$ gradus revera $m$ radices habere.

Probe itaque iis, quae hucusque exposita sunt, perpensis, demonstrationem novam theorematis gravissimi ex principiis omnino diversis petitam peritis haud ingratam fore spero, quam exponere statim aggredior.

\Article{13}
\textsc{Lemma.} \textit{Denotante $m$ numerum integrum positivum quemcunque, functio $\sin\varphi\mathbin{.}x^m-\sin m\varphi\mathbin{.}r^{m-1}x+\sin(m-1)\varphi\mathbin{.}r^m$ divisibilis erit per $xx-2\cos\varphi\mathbin{.}rx+rr$.}

\textit{Demonstr.} Pro $m=1$ functio illa fit $=0$ adeoque per quemcunque factorem divisibilis; pro $m=2$ quotiens fit $\sin\varphi$, et pro quovis valore maiori quotiens erit $\sin\varphi\mathbin{.}x^{m-2}+\sin2\varphi\mathbin{.}rx^{m-3}+\sin3\varphi\mathbin{.}rrx^{m-4}+\text{etc.}+\sin(m-1)\varphi\mathbin{.}r^{m-2}$. Facile enim confirmatur, multiplicata hac functione per $xx-2\cos\varphi\mathbin{.}rx+rr$, productum functioni propositae aequale fieri.
